% Imporditud: tools/import_eksperimendid.py
% Allikas: Eesti füüsikaolümpiaadi eksperimendi ülesannete
%   kogu (Rael Kalda, 2023), ülesanne Ü162, lahendus L162.
\setAuthor{EFO žürii}
\setRound{lõppvoor}
\setYear{2002}
\setNumber{G E2}
\setDifficulty{5}
\setTopic{dünaamika}
\setVahendid{statiiv, joonlaud, pliiats niidiga, klaasplaat, lapp klaasplaadi puhastamiseks}
\setPunktid{14}
\setAllikas{Eksperimendi kogu Ü162, lahendus lk 121}
\setLahendusseis{toores}

\prob{Hõõrdetegur}
Määrata pliiatsi teritatud otsa ja klaasplaadi vaheline hõõrdetegur.
\textit{Märkus:} pliiatsi võib lugeda ühtlase massijaotusega pulgaks.

\hint

\solu
Niidi üks ots on juba seotud pliiatsi tagumise otsa külge, niidi teise otsa kinnitame statiivi külge. Niidi pikkuse l (statiivist pliiatsini) võtame sama pika kui pliiats ise. Laseme statiivi tasapisi allapoole: alguses puudutab pliiats klaasplaati otse kinnituspunkti all (edasise töö mugavuse huvides võime selle punkti ära märkida asetades klaasplaadi alla paberi ristikesega), edasi hakkab ta viltu kalduma. Lõpuks hakkab pliiats libisema: fikseerime selle asendi ning mõõdame kauguse plaadist statiivi kinnituspunktini h. Kirjutades välja jõumomentide tasakaalu niidi keskpunkti suhtes näeme, et toereaktsiooni ja hõõrdejõu resultant läheb läbi niidi keskpunkti. Seega on hõõrdetegur selle nurga tangens, mis jääb pliiatsi toetuspunktist

tõmmatud vertikaali ja toetuspunkti ning niidi keskpunkti õhendava sirge vahele. Selle geomeetria ülesande lahendamisel leiame, et

1 $l_2$ $\mu$= 4 $h_2$ $-$ 1.

Aktsepteeritav on ka, kui valemi asemel tehakse sobivas mõõtkavas täpne joonis, mille pealt mõõdetakse otsitava nurga tangensi vahetuks arvutamiseks vajalikud vahemaad. Mõistlik tulemus on $\mu \approx$ 0,15.

Alternatiivlahendused: (a) Sama, mis eelmine, aga kinnituspunkti alla lastes lühendame ka niiti, nii et niit oleks kogu aeg horisontaalne (horisontaalsust võib kontrollida joonlaua abil). $\mu$ = l/(2L), kus L on pliiatsi pikkus. (b) h võtame võrdseks pliiatsi pikkusega L ning pikendame niiti libisemiseni.

$\mu$= l .

4L2 $-$ $l_2$
\probend
